\chapter{Related works and tools}

Chương này phân tích các nền tảng công nghệ và công cụ mã nguồn mở được tích hợp trong hệ thống đề xuất. Việc lựa chọn các công cụ này dựa trên các tiêu chí về khả năng mở rộng, tính riêng tư dữ liệu và hiệu suất thực thi trong môi trường cục bộ.

\section{Công cụ Đánh giá Tư thế Bảo mật Đám mây (CSPM)}

\subsection{Prowler}
Prowler là một công cụ dòng lệnh (CLI) mã nguồn mở tiêu chuẩn công nghiệp dùng để đánh giá bảo mật AWS. Nó được xây dựng dựa trên Khung kiến trúc tối ưu của AWS (AWS Well-Architected Framework) và các tiêu chuẩn CIS (Center for Internet Security).

Trong kiến trúc hệ thống Multi-Agent của đồ án, Prowler đóng vai trò là \textbf{Cảm biến (Sensor)}, cung cấp dữ liệu đầu vào thô cho quá trình nhận thức của tác nhân. Khác với các công cụ như AWS Security Hub (tốn chi phí) hay ScoutSuite (tập trung vào visualization), Prowler có khả năng xuất dữ liệu dưới định dạng \textbf{OCSF (Open Cybersecurity Schema Framework)}.

Định dạng OCSF giúp chuẩn hóa các phát hiện bảo mật (findings) thành cấu trúc JSON thống nhất, cho phép \textit{Monitoring Agent} và \textit{Remediate Agent} dễ dàng trích xuất các trường quan trọng như \texttt{event\_code}, \texttt{resource\_id}, và \texttt{status\_code} mà không cần xử lý ngôn ngữ tự nhiên phức tạp.

\begin{figure}[H]
  \centering
  \framebox{\parbox{0.8\textwidth}{\centering
    \vspace{1.5cm}
    \textbf{[Sơ đồ: Prowler Architecture]} \\
    \small\textit{Prowler Core $\rightarrow$ AWS API Calls $\rightarrow$ Compliance Checks $\rightarrow$ JSON-OCSF Output}
    \vspace{1.5cm}
  }}
  \caption{Luồng dữ liệu kiểm tra tuân thủ của Prowler}
  \label{fig:prowler_arch}
\end{figure}

\section{Hạ tầng Suy luận Mô hình Ngôn ngữ Cục bộ}

\subsection{Ollama và Llama 3.2}
Để đảm bảo tính riêng tư dữ liệu—một yêu cầu tối thượng trong kiểm toán bảo mật—hệ thống sử dụng \textbf{Ollama} làm nền tảng suy luận cục bộ (Local Inference Runtime). Ollama giải quyết bài toán phức tạp của việc quản lý bộ nhớ và tải mô hình thông qua việc đóng gói các mô hình vào định dạng tiêu chuẩn.

Mô hình được lựa chọn là \textbf{Llama 3.2} (phiên bản 3B tham số). Đây là một Mô hình Ngôn ngữ Nhỏ (SLM) được tối ưu hóa thông qua kỹ thuật \textit{Knowledge Distillation} và \textit{Pruning} từ các mô hình lớn hơn (Llama 3.1 70B). 
\begin{itemize}
    \item \textbf{Cơ chế Attention:} Llama 3.2 sử dụng kiến trúc \textit{Grouped-Query Attention (GQA)}, giúp giảm đáng kể băng thông bộ nhớ (VRAM) cần thiết trong quá trình giải mã (decoding), cho phép tốc độ sinh token nhanh hơn trên phần cứng giới hạn.
    \item \textbf{Instruction Tuning:} Mô hình đã được tinh chỉnh để tuân thủ các chỉ dẫn phức tạp, phù hợp cho việc phân tích ý định người dùng trong \textit{Planning Agent} và tạo kế hoạch JSON trong \textit{Remediate Agent}.
\end{itemize}

\section{Quản lý Dữ liệu và Bộ nhớ (Data Serialization)}
Hệ thống không sử dụng Database quan hệ hay Vector DB mà áp dụng cơ chế lưu trữ dựa trên tệp (File-based Storage) với định dạng **JSON**, phù hợp với kiến trúc Micro-services cục bộ:

\begin{itemize}
    \item \textbf{Tính cấu trúc:} JSON (JavaScript Object Notation) tương thích hoàn toàn với đầu ra của Prowler và đầu vào của LLM (thông qua LangChain Output Parsers).
    \item \textbf{Lưu vết (State Persistence):} Các tệp dữ liệu đóng vai trò là "Snapshots" của hệ thống tại từng thời điểm:
    \begin{itemize}
        \item \texttt{pre\_scan.json}: Trạng thái hạ tầng trước khi can thiệp.
        \item \texttt{analysis\_diff.json}: Bản ghi sự khác biệt (Delta) dùng để sinh báo cáo.
    \end{itemize}
    \item \textbf{Hiệu năng:} Với quy mô quét cụ thể (Scoped Scan), việc đọc/ghi file JSON trực tiếp trên ổ cứng (I/O) nhanh hơn và ít tốn tài nguyên hơn đáng kể so với việc duy trì kết nối đến một Database Server.
\end{itemize}

\section{AWS SDK cho Tự động hóa (Boto3)}
Boto3 là bộ công cụ phát triển phần mềm (SDK) chính thức của AWS dành cho Python. Trong hệ thống này, Boto3 đóng vai trò là \textbf{Cơ cấu chấp hành (Actuator)}.

Các tác nhân AI không trực tiếp thay đổi môi trường đám mây. Thay vào đó, \textit{Remediate Agent} tạo ra các tham số cấu hình, và các hàm \textit{Tool} (được viết bằng Boto3) sẽ thực thi các thay đổi đó một cách an toàn. Boto3 cung cấp các cơ chế xử lý lỗi (Error Handling) và quản lý phiên (Session Management) ở mức thấp, giúp đảm bảo tính ổn định khi thực hiện các tác vụ nhạy cảm như \texttt{put\_bucket\_versioning} hay \texttt{put\_user\_policy}.


